
\section{Experimental Details}
\label{app:experiment-details}

\subsection{Shared setup}
\label{app:shared-setup}

\paragraph{Variational family.}
All SIVI variants use a conditional Gaussian semi-implicit variational family.  We write the mixing variable as $\epsilon$ and the latent variable as $z$, with
\begin{equation}
  q_\phi(z)
  =
  \int q_\phi(z \mid \epsilon) q(\epsilon)\,d\epsilon,
  \qquad
  q_\phi(z \mid \epsilon)
  =
  \mathcal{N}\!\left(z;\mu_\phi(\epsilon),\Sigma_\phi\right).
\end{equation}
For two-dimensional targets, $\epsilon \in \mathbb{R}^{4}$ is drawn from a standard Gaussian prior, and $\mu_\phi$ and the diagonal variance are produced by a two-hidden-layer MLP with hidden width 128.  For the conditioned diffusion task, $\epsilon \in [0,1]^{100}$ uses a uniform prior, $\mu_\phi$ is a two-hidden-layer ReLU MLP with hidden width 128, and the conditional covariance is a learned global diagonal variance initialized near $0.01725$.  For BNN regression, $\epsilon \in \mathbb{R}^{3}$ uses a standard Gaussian prior, $\mu_\phi$ and the diagonal log variance are produced by a two-hidden-layer ReLU MLP with hidden width 10, and the log variance is initialized at $-8$ with a lower clamp at $-20$.  These BNN choices follow the compact semi-implicit architecture used in the KSIVI and SIVI-SM comparisons.

\paragraph{Common training protocol.}
All reported runs use ten random seeds; tables report the mean and standard error.  Unless otherwise noted, SIVI, UIVI, AISIVI, and DIVI use 10,000 main VI iterations with batch size 128, Adam optimization, and a step learning-rate scheduler. And unless otherwise noted, all methods are trained with annealing enabled, for half of the total training iterations. All experiments are conducted on a Linux server with Nvidia RTX 4090 GPUs and Intel Xeon Gold 6430 @ 2.10GHz.

\paragraph{SIVI.}
SIVI approximates the semi-implicit marginal $\log q_\phi(z)$ using Monte Carlo averages over auxiliary $\epsilon$ samples.  We draw 4,096 auxiliary samples on diffusion targets and 2,048 on BNN targets due to GPU memory constraints.

\paragraph{UIVI.}
UIVI samples from the reverse conditional $q_\phi(\epsilon \mid z)$ via Hamiltonian Monte Carlo.  After 5 burn-in transitions, a single HMC trajectory of 5 leapfrog steps provides $K=5$ auxiliary samples, one per intermediate position, following \citet{titsias2019unbiased}.  The leapfrog step size is 0.2.

\subsection{Metrics}
\label{app:metric-details}

\paragraph{KL-style ELBO.}
For toy and diffusion targets, we report a KL-style objective,
\begin{equation}
  D_{\mathrm{KL}}
  =
  -\widehat{\mathcal{L}}(q_\phi)
  =
  -\frac{1}{N}
  \sum_{i=1}^{N}
  \left[
    \log p(z_i) - \widehat{\log q_\phi(z_i)}
  \right],
  \qquad
  z_i \sim q_\phi .
\end{equation}
The estimate uses $N=5{,}000$ VI samples.  For each $z_i$, $\widehat{\log q_\phi(z_i)}$ is approximated with 20 batches of 2,048 auxiliary $\epsilon$ samples, plus the generating $\epsilon$ sample. We utilize the batching logic in \cite{pielok2025revisiting} to increase the effective number of Monte Carlo auxiliary samples to 40960.

\paragraph{Sliced Wasserstein-2.}
Sliced Wasserstein-2 is computed by averaging one-dimensional Wasserstein distances over random unit projections,
\begin{equation}
  \mathrm{SW}_2^2(q,r)
  =
  \mathbb{E}_{\theta \sim \mathrm{Unif}(\mathbb{S}^{d-1})}
  \left[
    W_2^2\!\left(\theta^\top q,\theta^\top r\right)
  \right],
\end{equation}
estimated using 10,000 samples and 5,000 random projections.  For heavy-tailed distributions, we report a truncated variant following \cite{li2023samplingmollifiedinteractionenergy}: both variational and reference samples are retained only if every coordinate has absolute value less than a threshold, the \textit{edge length}, and sliced Wasserstein-2 is then computed on the accepted samples.  This provides a stable central-region transport diagnostic when tail draws dominate finite-sample Wasserstein estimates.  We use 6 as the edge length for \texttt{8\_gaussians} and edge length 8 for \texttt{student\_uc}. 

\paragraph{Expected log marginal (ELM).}
For the conditioned diffusion benchmark, we report the expected log marginal likelihood,
\begin{equation}
  \mathrm{ELM}
  =
  \mathbb{E}_{z \sim r}
  \left[
    \sum_{j=1}^{d}
    \log \widehat{q}_{\phi,j}(z_j)
  \right],
\end{equation}
where $r$ is the empirical distribution from the SGLD reference samples and $\widehat{q}_{\phi,j}$ is a one-dimensional Gaussian KDE fit to coordinate $j$ of model samples.  We follow \citet{pielok2025revisiting} in using this metric.  Direct Monte Carlo approximation of the full semi-implicit marginal likelihood is difficult because it integrates over a 100-dimensional mixing distribution, so the coordinate-wise KDE estimate provides a stable diagnostic, though it does not capture the full joint marginal likelihood. The reported evaluation uses 1,000 reference samples and 100,000 model samples; the SGLD baseline uses 100,000 SGLD samples as the model sample set.

\paragraph{BNN predictive metrics.}
For BNN targets, evaluation draws 2,000 posterior samples (network weights) from the learned variational distribution.  RMSE is computed from the posterior predictive mean on the test set.  Test log-likelihood is estimated by Monte Carlo marginalization over posterior samples, and the reported NLL is its negative.

\subsection{Two-dimensional targets}
\label{app:toy-details}

The main text displays and summarizes \texttt{x\_shaped}, \texttt{student\_uc}, and \texttt{8\_gaussians}.  Visual panels use 2,000 samples per method and overlay the samples on target-density contours.

\paragraph{Target density definitions.}
We now state the precise density for each of the three two-dimensional synthetic targets used in Section~\ref{subsec:toy-experiments}.
All targets have $z\in\mathbb{R}^2$.

\subparagraph{X-shaped.}
The \texttt{x\_shaped} target is an equal-weight mixture of two zero-mean Gaussians whose covariance ellipses cross at the origin, producing an X-shaped density:
\begin{equation}\label{eq:x-shaped}
  p(z) = \tfrac{1}{2}\,\mathcal{N}(z;\,0,\,\Sigma_1)
        + \tfrac{1}{2}\,\mathcal{N}(z;\,0,\,\Sigma_2),
\end{equation}
where
\[
  \Sigma_1 = \begin{pmatrix} 2 & 1.8 \\ 1.8 & 2 \end{pmatrix},
  \qquad
  \Sigma_2 = \begin{pmatrix} 2 & -1.8 \\ -1.8 & 2 \end{pmatrix}.
\]
The high off-diagonal correlation ($\rho = \pm 0.9$) in each component creates two narrow ridges that intersect at the origin.  Approximating both ridges simultaneously requires the variational family to represent sharp, non-convex geometry.

\subparagraph{Student-t (uncorrelated marginals).}
The \texttt{student\_uc} target is a bivariate Student-$t$ distribution with $\nu=2$ degrees of freedom and a fixed scale matrix $A\in\mathbb{R}^{2\times 2}$ with $\det(A)=1$.  Its density is defined through the transformation $\tilde{z} = A^{-1}z$ with independent univariate Student-$t$ marginals:
\begin{equation}\label{eq:student-uc}
  p(z) \;=\; |\!\det A|^{-1}\,\prod_{i=1}^{2}
    \frac{\Gamma\!\bigl(\tfrac{\nu+1}{2}\bigr)}
         {\sqrt{\nu\pi}\;\Gamma\!\bigl(\tfrac{\nu}{2}\bigr)}
    \left(1 + \frac{\tilde{z}_i^{\,2}}{\nu}\right)^{\!-(\nu+1)/2},
  \qquad \nu = 2.
\end{equation}
The matrix $A$ is generated deterministically (seed~50) as $A = (B B^\top)^{-1}/\det\bigl((B B^\top)^{-1}\bigr)$, where $B_{ij}\sim\mathrm{Uniform}(-1,1)$.  The resulting density has polynomial tails ($\propto\|z\|^{-3}$ marginally), which stress-tests whether the variational approximation can allocate sufficient probability mass far from the mode.

\subparagraph{Eight Gaussians.}
The \texttt{8\_gaussians} target is an equally weighted mixture of eight isotropic Gaussians arranged uniformly on a circle of radius~$r$:
\begin{equation}\label{eq:8-gaussians}
  p(z) = \frac{1}{8}\sum_{k=1}^{8}
    \mathcal{N}\!\bigl(z;\,\mu_k,\,\sigma^2 I_2\bigr),
  \qquad
  \mu_k = r\bigl(\cos(k\pi/4),\;\sin(k\pi/4)\bigr),
\end{equation}
with $r=4$ and $\sigma=0.5$.  The well-separated modes (inter-mode distance $\approx 3.1$, versus component standard deviation $0.5$) make this a challenging multi-modal target: the variational family must discover and cover all eight modes without spurious bridging density between them.

\bigskip

Table~\ref{tab:toy-target-summary} collects the key parameters of each target.

\begin{table}[h]
\centering
\caption{Summary of two-dimensional synthetic target densities.}
\label{tab:toy-target-summary}
\small
\begin{tabular}{@{}llll@{}}
\toprule
\textbf{Target} & \textbf{Density form} & \textbf{Parameters} & \textbf{Key challenge} \\
\midrule
\texttt{x\_shaped}
  & $\frac{1}{2}\mathcal{N}(0,\Sigma_1) + \frac{1}{2}\mathcal{N}(0,\Sigma_2)$
  & $\Sigma_{1,2} = \bigl(\begin{smallmatrix}2 & \pm1.8\\ \pm1.8 & 2\end{smallmatrix}\bigr)$
  & Intersecting ridges \\[6pt]
\texttt{student\_uc}
  & $|A|^{-1}\prod_{i=1}^{2} t_\nu(\tilde{z}_i)$, \; $\tilde{z}=A^{-1}z$
  & $\nu=2$, \; $\det(A)=1$
  & Heavy polynomial tails \\[6pt]
\texttt{8\_gaussians}
  & $\frac{1}{8}\sum_{k=1}^{8}\mathcal{N}(\mu_k, \sigma^2 I)$
  & $r=4$, \; $\sigma=0.5$
  & Multi-modal coverage \\
\bottomrule
\end{tabular}
\end{table}


\paragraph{Training configuration details.}
All methods use 10,000 iterations, VI learning rate $10^{-3}$, batch size 128, and a step scheduler with decay factor 0.7 every 1,000 iterations (0.9 on \texttt{8\_gaussians}).  Linear annealing over 5,000 steps is applied on all targets except \texttt{student\_uc}, where annealing is disabled for all methods.

\subparagraph{AISIVI.}
AISIVI uses a conditional RealNVP reverse model with 8 coupling layers and hidden width 12.  Before joint training, the flow is pre-trained on initial VI samples for 500 warmup epochs with learning rate $10^{-2}$ and batch size 8,192.  During main training, the flow uses 1,024 reverse samples, reverse learning rate $10^{-3}$, and one reverse update per VI step.

\subparagraph{DIVI.}
The proxy score MLP has hidden width 64 and three hidden layers.  The proxy is warm-started for 500 epochs with learning rate $10^{-3}$ and batch size 8,192.  During main training the proxy uses reverse batch size 2,048, reverse learning rate $10^{-3}$, and 5 updates per VI step.

% \paragraph{KSIVI.}
% KSIVI is trained for 50,000 iterations with 25,000-step linear annealing (except on \texttt{student\_uc}), a step-decay schedule with factor 0.9 every 1,000 iterations, and the default settings from \citet{cheng2024kernel}.  We use a Gaussian kernel on most targets, an IMQ kernel on \texttt{multimodal}, and a Riesz kernel on \texttt{student\_uc}.

\paragraph{Quantitative results.}
We here present the quantitative results in Table~\ref{tab:toy-method-grid}. It can be seen from the table that, DIVI obtains the best metric on all targets in terms of KL divergence and Wasserstein-2 distance. Here the W2 distance evaluation follows the protocol in Appendix~\ref{app:metric-details}. We are observing divergence on \texttt{8_gaussians} for AISIVI, as a result of drifting support. The same divergence behavior has been observed in the official code of AISIVI, indicating a potential weakness in training stability.  

\begin{table}[t]
\centering
\caption{Selected toy target final evaluation metrics by method.}
\label{tab:toy-method-grid}
\begin{tabular}{lccccccc}
\toprule
Method & \multicolumn{3}{c}{ELBO} & \multicolumn{3}{c}{W2} & Avg. training time (s) \\
 & student_uc & 8_gaussians & x_shaped & student_uc & 8_gaussians & x_shaped &  \\
\midrule
UIVI & -2.683 $\pm$ 0.00 & -0.804 $\pm$ 0.01 & -0.063 $\pm$ 0.03 & 0.141 $\pm$ 0.01 & 0.328 $\pm$ 0.05 & 0.207 $\pm$ 0.06 & 752.539 $\pm$ 6.23 \\
AISIVI & -2.681 $\pm$ 0.00 & -17.685 $\pm$ 10.89 & -0.076 $\pm$ 0.03 & 0.125 $\pm$ 0.01 & 0.672 $\pm$ 0.17 & 0.196 $\pm$ 0.06 & 335.754 $\pm$ 5.53 \\
DSIVI & \textbf{-2.677 $\pm$ 0.00} & \textbf{-0.389 $\pm$ 0.15} & \textbf{-0.008 $\pm$ 7.99e-04} & \textbf{0.083 $\pm$ 0.00} & \textbf{0.318 $\pm$ 0.11} & \textbf{0.075 $\pm$ 0.01} & 182.031 $\pm$ 1.10 \\
\bottomrule
\end{tabular}
\end{table}



\subsection{Conditioned diffusion process}
\label{app:diffusion-details}

The diffusion posterior is the 100-dimensional discretized path posterior described in Section~\ref{subsec:langevin-experiments}.  Reference samples are produced by SGLD with 1,000 particles, 100,000 iterations, and step size $10^{-4}$.  Trace plots show posterior mean paths and marginal 95\% posterior intervals against the true latent path and noisy observations.

All reverse-model methods (SIVI, UIVI, AISIVI, DIVI) are trained for 10,000 VI iterations with learning rate $10^{-3}$, batch size 128, step scheduler (decay 0.7 every 1,000 steps), and 5,000-step linear annealing.  KSIVI is trained for 100,000 iterations.

\paragraph{AISIVI.}
AISIVI uses a conditional RealNVP with 16 coupling layers and hidden width 256.  Before joint training, the reverse flow is pre-trained on samples from the initial variational distribution for 2,000 warmup epochs with learning rate $2 \times 10^{-4}$ and batch size 128.  During main training, the flow uses 256 reverse samples, reverse learning rate $2 \times 10^{-4}$, and one reverse update per VI step.

\paragraph{DIVI.}
The proxy score MLP has hidden width 128 and three hidden layers.  The proxy is warm-started for 2,000 epochs with learning rate $10^{-3}$ and batch size 8,192.  During main training the proxy uses reverse batch size 2,048, reverse learning rate $10^{-3}$, and 5 updates per VI step.

\paragraph{KSIVI.}
KSIVI follows the default schedule from \citet{cheng2024kernel}: 100,000 iterations, 50,000-step linear annealing, Gaussian kernel, and a step-decay schedule with factor 0.9 every 10,000 iterations.

\paragraph{Evaluation.}
We evaluate using the expected log marginal likelihood defined in Appendix~\ref{app:metric-details}, with 1,000 reference samples from the SGLD chain and 100,000 model samples from each trained variational distribution.

\subsection{Bayesian neural-network regression}
\label{app:bnn-details}

Each BNN target uses a one-hidden-layer ReLU network with 50 hidden units.  The input dimensions are 13 for Boston, 8 for Concrete, 4 for Power, 9 for Protein, 11 for Wine Red, and 6 for Yacht. The datasets can be found on the UCI Machine Learning website \footnote{\url{https://archive.ics.uci.edu/}}, and all follow CC BY 4.0 License to the best of our knowledge. Including weights, biases, and two precision variables, the corresponding posterior dimensions are 751, 501, 301, 551, 651, and 401.  The datasets contain 506, 1,030, 9,568, 45,730, 1,599, and 308 observations respectively.  Data preparation uses a fixed random seed, a 90/10 train-test split, input standardization based on the training split, and target standardization for training labels.  Test labels are evaluated on their original scale.

\paragraph{Common BNN protocol.}
All methods use batch size 128, gradient clipping at norm 1, and a 100-step MSE pretraining stage that initializes the variational sampler toward predictive weights before posterior training begins.  For evaluation, we apply exponential moving averages of the VI parameters with decay 0.999, following the stabilized BNN evaluation protocol used by KSIVI and SIVI-SM \citep{cheng2024kernel,yu2023semi}.  The VI learning rate is selected separately for each method and target from $\{10^{-3}, 10^{-4}, 10^{-5}\}$ by cross-validation, since high-learning-rate divergence depends on both the method and the target posterior.

\paragraph{AISIVI.}
AISIVI uses a conditional RealNVP with 16 coupling layers and hidden width 256.  The reverse flow is pre-trained on initial VI samples for 500 warmup epochs with learning rate $5 \times 10^{-5}$ and batch size 100.  During main training, the flow uses 256 reverse samples, reverse learning rate $5 \times 10^{-5}$, one reverse update per VI step, and 10,000-step linear annealing.

\paragraph{DIVI.}
The proxy score MLP has hidden width 512 and four hidden layers.  The proxy is warm-started for 5,000 epochs with learning rate $2 \times 10^{-3}$ and batch size 2,048.  During main training the proxy uses 2 updates per VI step.  Linear annealing over 10,000 steps is applied.

\paragraph{KSIVI.}
KSIVI is trained for 100,000 iterations without annealing, using a Gaussian kernel and log-likelihood regularization term with weight 1, with a step-decay schedule of factor 0.9 every 3,000 iterations, following \citet{cheng2024kernel}.

\paragraph{Evaluation.}
BNN evaluation draws 2,000 posterior samples and reports test RMSE and NLL as described in Appendix~\ref{app:metric-details}.

% \subsection{Reproducibility}
% \label{app:reproducibility-details}

% All results in the tables and figures are generated from the same evaluation workflow and seed set.  The standalone draft includes the generated figure and table inputs through the local LaTeX wrapper; the paper source should use the corresponding generated PDF figures and table inputs when integrated into the full manuscript.

\section{Numerical Validations}
\label{app: numerical-experiments}

In this section, we provide empirical diagnostics supporting Assumption~\ref{assump:bounded_score} and Assumption~\ref{assump:bounded_reparam} through numerical experiments along the observed training trajectories. 

For Assumption~\ref{assump:bounded_score}, we provide empirical diagnostics by tracking the fourth moment of scores, both for the reverse model and for the target score. We plot the mean value of the fourth moment of norms of the reverse model score $\mathbb{E}_{z\sim q_\phi} \lVert s_\psi(z) \rVert^4$ as the solid line, and the minimum and maximum values as the shaded regions in Figure~\ref{fig:score_q_4th_moment}, with each of the values calculated across 10 independent seeds. The same plot for fourth moment of the target score $\mathbb{E}_{z\sim q_\phi} \lVert \nabla_z \log p(\mathcal{D}, z) \rVert^4$ is shown in Figure~\ref{fig:score_p_4th_moment}. The figures provide empirical evidence that there exists a uniform bound on the $4^{\text{th}}$ moment of the norm of the scores over the training trajectories, consistent with  Assumption~\ref{assump:bounded_score}.

% \begin{figure}[t]
%     \centering
%     \includegraphics[width=0.8\linewidth]{weight_norm_iteration_grid.pdf}
%     \caption{Norm of the VI model parameters throughout DIVI training for \texttt{x\_shaped} and the conditioned diffusion benchmark, aggregated over 10 seeds.}
%     \label{fig:weight_norm}
% \end{figure}

\begin{figure}[t]
    \centering
    \includegraphics[width=0.8\linewidth]{score_q_4th_moment_iteration_grid.pdf}
    \caption{Fourth moment of the norm of the approximated score $\mathbb{E}_{z\sim q_\phi} \lVert s_\psi(z) \rVert^4$ throughout DIVI training for \texttt{x\_shaped} and the conditioned diffusion benchmark, aggregated over 10 seeds.}
    \label{fig:score_q_4th_moment}
\end{figure}

\begin{figure}[t]
    \centering
    \includegraphics[width=0.8\linewidth]{score_p_4th_moment_iteration_grid.pdf}
    \caption{Fourth moment of the norm of the target score $\mathbb{E}_{z\sim q_\phi} \lVert \nabla_z \log p(\mathcal{D}, z) \rVert^4$ throughout DIVI training for \texttt{x\_shaped} and the conditioned diffusion benchmark, aggregated over 10 seeds.}
    \label{fig:score_p_4th_moment}
\end{figure}


For Assumption~\ref{assump:bounded_reparam}, we report the calculated $\max \big\{ \mathbb{E}_{\epsilon} \lVert \nabla_\phi \sigma_\phi(\epsilon)\rVert_2^4, \mathbb{E}_\epsilon \lVert \nabla_\phi \mu_\phi(\epsilon) \rVert_2^4\big\}$ for each checkpoint respectively throughout the training process, and aggregate across 10 independent seeds. We track the metrics for the toy \texttt{x_shaped} and Conditioned Diffusion target distribution. The results are shown in Figure~\ref{fig:bounded_m_eps}, where the solid line indicates the mean while the shaded region demonstrates the minimum and maximum values. These diagnostics suggest that we can obtain a uniform upper bound for the expected $\ell_2$ norm's fourth moment of the VI networks' gradient, justifying Assumption~\ref{assump:bounded_reparam}.

\begin{figure}[t]
    \centering
    \includegraphics[width=0.8\linewidth]{m_eps_iteration_grid.pdf}
    \caption{The expected $\ell_2$ norm's fourth moment of the VI model gradient throughout the DIVI training process for \texttt{x_shaped} and Conditioned Diffusion experiments, aggregated across 10 different seeds. This supports Assumption~\ref{assump:bounded_reparam} that the VI model gradients are bounded. }
    \label{fig:bounded_m_eps}
\end{figure}

We also provide empirical evidence to the alternative version of Assumption~\ref{assump:bounded_score}, as present in Appendix~\ref{app:linear-growth}. In Appendix~\ref{app:linear-growth} we proved that Assumption~\ref{assump:bounded_score} can be derived from another linear growth condition $\max\{\lVert \nabla_z \log p(\mathcal{D},z) \rVert, \lVert s_\psi(z)\rVert \} \leq \lambda (\lVert z \rVert + 1)$ together with a bounded fourth-moment condition $\max\{ \mathbb{E}_\epsilon \lVert \mu_\phi(\epsilon) \rVert^4, \mathbb{E}_\epsilon \lVert \sigma_\phi(\epsilon) \rVert^4 \} \leq \zeta $, both of which we assume to hold uniformly over $\phi$. As empirical evidence, we demonstrate for three checkpoints taken at 2000, 5000 and 10000 iterations, that $\log \max\{\lVert \nabla_z \log p(\mathcal{D},z) \rVert, \lVert s_\psi(z)\rVert \} - \log (\lVert z \rVert + 1)$ is bounded in Figure~\ref{fig:score_norm_linearity}. The figure demonstrates different values for samples of $z$ drawn with uniform distribution in $\lVert z \rVert$.  

We also plot the $4^{\text{th}}$ moment of the VI model checkpoints $\max\{ \mathbb{E}_\epsilon \lVert \mu_\phi(\epsilon) \rVert^4, \mathbb{E}_\epsilon \lVert \sigma_\phi(\epsilon) \rVert^4 \}$ using 10 different seeds and plot them along the 10 checkpoints we obtained during the DIVI training regime in Figure~\ref{fig:vi_fourth_moment}. The solid line indicates the 10-seed-average of the fourth moments, while the shaded regions indicate the minimum and maximum of recorded values. These plots together demonstrate empirical evidence to support the alternate version of Assumption~\ref{assump:bounded_score} in Appendix~\ref{app:linear-growth}.


\begin{figure}[t]
    \centering
    \includegraphics[width=0.8\linewidth]{score_norm_linearity_uniform_grid.pdf}
    \caption{$\log \max\{\lVert \nabla_z \log p(\mathcal{D},z) \rVert, \lVert s_\psi(z)\rVert \} - \log (\lVert z \rVert + 1)$ over different $z$'s, $\phi$'s and $\psi$'s throughout DIVI training for \texttt{x\_shaped} and the conditioned diffusion benchmark.}
    \label{fig:score_norm_linearity}
\end{figure}

\begin{figure}[t]
    \centering
    \includegraphics[width=0.8\linewidth]{vi_fourth_moment_iteration_grid.pdf}
    \caption{Fourth moment of the VI model reparameterization norm $\max\{ \mathbb{E}_\epsilon \lVert \mu_\phi(\epsilon) \rVert^4, \mathbb{E}_\epsilon \lVert \sigma_\phi(\epsilon) \rVert^4 \}$ throughout DIVI training for \texttt{x\_shaped} and the conditioned diffusion benchmark, aggregated over 10 seeds.}
    \label{fig:vi_fourth_moment}
\end{figure}
